\documentclass[11pt]{article}
% Shared preamble for private TrustAndReview manuscript drafts.
\usepackage[margin=1in]{geometry}
\usepackage[T1]{fontenc}
\usepackage[utf8]{inputenc}
\usepackage{lmodern}
\usepackage{microtype}
\usepackage{amsmath,amssymb,mathtools}
\usepackage{booktabs,tabularx,longtable}
\usepackage{enumitem}
\usepackage{xcolor}
\usepackage{hyperref}
\usepackage[round,authoryear]{natbib}
\hypersetup{colorlinks=true,linkcolor=blue!50!black,citecolor=blue!50!black,urlcolor=blue!60!black}
\setlist{nosep}

\newcommand{\draftstatus}[1]{%
  \begin{center}\fcolorbox{black}{yellow!15}{\parbox{0.92\linewidth}{\textbf{Private draft status:} #1}}\end{center}}
\newenvironment{evidencegate}[1][Evidence required before promotion]
  {\begin{quote}\color{red!60!black}\textbf{#1.} }
  {\end{quote}}
\newcommand{\designchoice}[1]{\textbf{Design choice:} #1}
\newcommand{\openhypothesis}[1]{\textbf{Open hypothesis:} #1}
\newcommand{\implementationobservation}[1]{\textbf{Implementation observation:} #1}


\usepackage{array}
\usepackage{caption}
\usepackage{float}
\usepackage{graphicx}
\usepackage{url}

\newcommand{\tdrg}{TDRG}
\newcommand{\mcrp}{MCRP}
\newcommand{\na}{\textsc{not available}}

\title{Plural Review Governance in Practice:\\
A Preregistered, Safety-Gated Pilot of Independent Trust Domains}
\author{Private working draft---human authorship and contribution record pending}
\date{24 August 2026}

\begin{document}
\maketitle

\draftstatus{Private Stage-0 protocol first draft. No ethics/IRB determination,
recruitment authorization, participant enrollment, intervention, data collection,
field result, or release approval is represented here. Until every applicable gate
is recorded, work is limited to non-participant protocol development and synthetic
fixtures. This document must not be used to recruit, deploy, or claim that the
intervention is useful, safe, fair, anonymous, or effective.}

\begin{abstract}
Machine-verifiable scientific records can make claims, evidence, provenance, and
review events inspectable, but protocol conformance alone cannot establish how a
review system behaves in scientific communities shaped by expertise, hierarchy,
scarcity, privacy risk, and conflicting institutional authority. We specify a
preregisterable, safety-gated sociotechnical pilot of Trust-Domain and Review
Governance (\tdrg{}), a proposed arrangement in which independently governed
communities interpret a shared claim/evidence record through local, topic- and
time-scoped trust overlays rather than a universal reviewer score. The first human
stage, if separately approved, would be a consented shadow pilot: experimental
routing, private standing, and review explanations would run alongside ordinary
review but could not determine the official scientific disposition. The protocol
defines independent-actor and management-control audits, consent and anti-coercion
protections, an adversary-specific anonymity threat model, appeals, adverse-event
and stopping rules, mixed-method outcomes, multiaxis resource accounting, and a
disclosure-safe reproducibility package. Co-primary outcomes, allocation, sample
size, participating domains, instruments, privacy thresholds, and numerical stop
bounds remain intentionally unset pending ethics, methods, privacy, and community
review. No participant results exist. The proposed pilot is designed to decide
whether measurement is credible and whether a larger operational study could be
ethical---not to certify \tdrg{} as superior or to rank reviewers as people.
\end{abstract}

\section{Introduction}

Scientific review is simultaneously epistemic work and organized labor. A review
may challenge a claim, inspect a workflow, identify a missing calibration, or trace
a transitive tool; it is also assigned, interpreted, and acted upon by institutions
whose members differ in authority, security, time, and freedom to dissent. A record
that perfectly preserves a review does not establish that the reviewer was
competent, independent, or safe from retaliation. Conversely, an organizational
decision about reviewer standing does not establish whether the reviewed claim is
true.

Interventions in peer-review procedure can change different outcomes in different
directions. In a randomized BMJ study, asking reviewers to be identified to authors
did not materially change the measured review-quality score or recommendation but
increased the likelihood that invitees declined to review
\citep{vanrooyen1999open}. In a controlled conference experiment, reviewers who
could see author identities exhibited different bidding and recommendation patterns
associated with author and institutional prestige \citep{tomkins2017bias}. These
studies do not predict the effect of \tdrg{}, but they demonstrate why acceptance,
review quality, participation, delay, and burden cannot be treated as one outcome.

The proposed intervention separates a shared, version-addressed claim/evidence and
review-event record from plural local views. Each view is rooted in a particular
community and conditioned on domain, topic, capability, role, case, policy, a
frozen event cut, and a separately declared evaluation clock. Communities may
import another domain's attestations under an explicit
bridge, reject them, or fork their policy without rewriting the shared record. A
reviewer uses a fresh case-scoped persona where the approved deployment supports
it; scientific feedback, private standing events, official dispositions, and
conduct or safety events remain distinct event types.

The Minimum Credible Reproducibility Protocol (MCRP) core governs the claim/evidence
scientific record. The Trust-Domain Review Graph (TDRG) is an optional, separately
versioned actor-routing and reliance profile; it remains outside MCRP core
conformance until its privacy, accountability, interoperability, replay, and
empirical promotion gates pass. Papers~01--03 own the design map, formal model, and
normative profile respectively; this paper owns only the ethics-gated
independent-actor pilot protocol.

\designchoice{The intervention has no universal actor score. The pilot will not
estimate whether a participant is a ``good reviewer'' or a trustworthy person. It
will estimate predeclared process outcomes for cases, assignments, episodes, and
domains under named policies and trust boundaries.}

This paper is a Stage-0 protocol and preregistration scaffold. It makes four bounded
contributions. First, it defines questions and estimands for a safety-first shadow
evaluation. Second, it operationalizes independence as a control graph rather than
an institutional label. Third, it specifies protections against coercion,
management capture, identity exposure, and automatic sanction. Fourth, it defines
how evidence may later be promoted into empirical claims without exposing
participants. It reports no observations about field performance.

\begin{evidencegate}[Human-subject gate]
No interview, survey, recorded tabletop exercise, shadow review, participant-linked
log, or operational trust decision may be collected for research until the
responsible ethics/IRB authority has issued a written determination and the approved
protocol, consent or waiver basis, privacy review, data plan, incident path,
compensation, and preregistration are recorded. The investigators may not decide
for themselves that the activity is exempt or outside human-subject regulation.
\end{evidencegate}

\section{Background and design rationale}

\subsection{Peer review has multiple, non-interchangeable outcomes}

Peer-review research does not supply a universal scalar definition of quality.
Open-review, blinding, training, assignment, and editorial interventions affect
participation, recommendations, report content, and timeliness through different
mechanisms. Systematic mapping of open peer review likewise describes a family of
practices rather than one intervention \citep{rosshellauer2017open}. The pilot
therefore avoids treating editor agreement, author satisfaction, majority vote,
citations, rank, or institutional prestige as ground truth. Material defects and
evidence gaps will be evaluated by a separate, uncertainty-preserving adjudication
procedure; process legitimacy, privacy, burden, and harms remain separate outcomes.

\subsection{Precommitment is necessary but not sufficient}

Preregistration and Registered Reports distinguish planned tests from analyses
chosen after outcomes are known \citep{nosek2018preregistration}. A blinded
comparison of published Registered Reports with comparison papers found higher
ratings on several methodological criteria, but the evidence arose from a bounded
sample and does not establish that preregistration prevents all bias or harm
\citep{soderberg2021registered}. Here preregistration serves narrower purposes: it
fixes the intervention, outcome hierarchy, disclosure rules, decision thresholds,
and amendment process before participant data are observed. Exploratory analyses
will remain permissible only when clearly labeled and separated from confirmatory
claims.

\subsection{Institutional power is part of the causal system}

Organizational-silence theory identifies management beliefs and structures as
conditions that can make speaking up appear unsafe \citep{morrison2000silence}.
Empirical participatory-design scholarship likewise treats participation and
decision authority as distinct: attendance does not imply control
\citep{bratteteig2014participation}. For this pilot, line management, employment,
funding, compensation, operator custody, agenda control, and authorship are not
``context'' to be adjusted away. They are measured pathways that may shape
voluntariness, candor, assignment, appeal, withdrawal, and interpretation.

\subsection{Plural governance is a hypothesis, not an endorsement}

Polycentric governance describes multiple decision centers operating with some
degree of autonomy under shared rules \citep{ostrom2010polycentric}. It motivates
the distinction between a shared event record and local governance views, but does
not prove that such pluralism improves peer review. Forking may localize capture or
preserve exit; it may also duplicate labor, fragment expertise, strengthen bridge
brokers, or permit accountability shopping. The pilot must preserve both
interpretations.

\subsection{Confidentiality and anonymity require named adversaries}

Privacy is an outcome of data processing and organizational practice, not a field
in a schema. The NIST Privacy Framework provides a risk- and outcome-based approach
to identifying problematic data actions and their consequences
\citep{nist2020privacy}. Reidentification research demonstrates that apparently
anonymous high-dimensional records may be linkable with auxiliary information
\citep{narayanan2008deanonymizing}; stylometric work shows that text itself can
carry identifying signals \citep{brennan2012stylometry}. Accordingly, every privacy
claim in this pilot is conditional on an adversary, accessible data, timing,
operator and collusion assumptions, attack method, observation window, and residual
risk. When a preregistered anonymity threshold is not met, the artifact will be
called \emph{confidential}, not anonymous.

\section{Research objectives and estimands}

The overarching question is: when independently governed scientific communities
use a shared claim/evidence record with local trust overlays, what changes in
review process, capture localization, minority and newcomer treatment, participant
burden, privacy, and institutional power compared with their ordinary review
process? The answer is unknown. The pilot is intended to determine feasibility,
safety, measurement quality, and mechanism plausibility.

The preregistration will nominate no more than three co-primary outcomes after
instrument testing and stakeholder review. Candidate research questions are:

\begin{enumerate}
  \item Does shadow \tdrg{} change adjudicated material-defect or evidence-gap
  yield per reviewer-hour relative to the ordinary process?
  \item Do independently governed domains derive different, inspectable routing
  or standing decisions from the same event record, and can participants explain
  those differences without seeing private histories?
  \item Does a staged governance fault or management-pressure scenario remain
  localized to the affected domain under the declared control graph?
  \item Are qualified dissenting reviews retained and discoverable without an
  automatic adverse standing update while the scientific dispute is unresolved?
  \item What waiting time, supervision load, false exclusion, progression, and
  exit patterns occur for newcomers and specialists in small pools?
  \item Under each declared adversary, what linkage, reidentification,
  duplicate-prevention, appeal, and authorized-opening outcomes occur?
  \item How do employment, funding, compensation, workload, and operator control
  affect reported candor, challenge, withdrawal, and exit?
  \item What human expertise, agent capability, governance labor, compute,
  storage, security, and support are required by each intervention function?
  \item What retaliation, chilling, exclusion, surveillance, gaming, false
  confidence, burden shifting, confidentiality, or security harms arise?
\end{enumerate}

\openhypothesis{A local-view architecture may contain some governance faults, but
shared operators, bridge authorities, funding, infrastructure, or personnel may
create cross-domain dependence. Localization is therefore an empirical outcome
conditioned on the audited control graph, not a property inferred from labels.}

Every outcome will name its unit and estimand. Candidate units are review cases,
review assignments, reviewers, governance episodes, and participating domains.
For example, a process estimand could be the domain-stratified difference in
adjudicated material findings per logged reviewer-hour between a shadow-\tdrg{}
assignment and the ordinary comparator during the same stage. A privacy estimand
could be adversary-specific linkage precision and recall over eligible case-persona
pairs. A burden estimand could be additional human and agent hours per completed
case, disaggregated by role. These examples do not fix the final estimands.

\begin{evidencegate}[Unresolved design parameters]
Participating domains, allocation unit, comparator details, co-primary outcomes,
smallest meaningful differences or precision targets, sample size, instruments,
privacy thresholds, and numerical stop bounds are unset. They must be chosen before
recruitment with independent statistics, qualitative-methods, domain, ethics, and
privacy review; placeholders must not be silently converted into analysis choices.
\end{evidencegate}

\section{Intervention, comparator, and independent actors}

\subsection{Frozen intervention package}

The intervention is the exact future release of the \tdrg{} protocol, evaluator,
schemas, policy overlays, identity/persona mechanism, and operator procedures
named in the deployment manifest. The manifest will include source and container
digests, configuration, randomness, clocks, external services, credential and
opening authorities, storage jurisdictions, and every transitive service admitted
inside the trust boundary. Any change is a versioned amendment; changes affecting
consent, threat assumptions, outcomes, or official disposition trigger review and
potential reconsent.

Four event types remain separate:
\begin{enumerate}
  \item scientific review observations about claims and evidence;
  \item private standing events used only inside an approved local view;
  \item official editorial or scientific dispositions made by the ordinary
  process; and
  \item conduct, safety, or policy events handled by the approved incident path.
\end{enumerate}
No raw vote, graph anomaly, scientific disagreement, or allegation automatically
creates a sanction or conduct finding.

\subsection{Comparator and shadow separation}

The comparator is each domain's documented ordinary review process, including its
assignment, conflict, adjudication, and appeal practices. ``Business as usual'' is
not a sufficient specification: the preregistration will enumerate comparator
steps and record concurrent changes. In Stage~1 the two processes may share a case
only under approved access rules, but experimental outputs cannot affect the
official disposition, employment, funding, publication, authorship, discipline, or
access to ordinary review.

Where ethical and operational review permits, allocation may occur by case,
assignment, or sequence and be stratified by domain, topic, case risk, and reviewer
experience. If randomization is infeasible, a matched, stepped-wedge,
interrupted-time, or other design may be used only with its identification
assumptions and causal limitations preregistered. No design will describe a
before/after association as a causal effect merely because the software changed.

\subsection{Control graph and domain independence}

An ``independent domain'' is not defined by a name, institution, or separate user
interface. Before deployment, each candidate domain must file a versioned control
manifest identifying:

\begin{itemize}
  \item governance roots, charter amendment, appointment, removal, and veto
  authority;
  \item line management, supervision, employment, funding, and compensation;
  \item credential, identity, assignment, hosting, data, model, opening, appeal,
  and incident operators;
  \item policy authorship, evaluator custody, infrastructure, and key management;
  \item overlapping personnel, conflicts, contractual obligations, and legal
  jurisdiction; and
  \item an exercisable withdrawal or fork path that preserves permissible
  scientific records without exporting private histories.
\end{itemize}

The manifests form a directed, typed control graph. Common controllers and shared
services are disclosed, not averaged into a nominal independence score. At least
two participating domains must have distinct governance roots and a credible exit
or fork path. Independence will be reported dimension by dimension and used in
sensitivity analysis. If a shared controller can alter both views, the affected
localization hypothesis is not testable as stated.

\section{Ethics, consent, positionality, and anti-coercion}

The Belmont principles connect respect for persons, beneficence, and justice to
informed consent, risk--benefit assessment, and participant selection
\citep{belmont1979}. The U.S. Common Rule supplies regulatory requirements where it
applies (45~CFR~46; official text at
\url{https://www.hhs.gov/ohrp/regulations-and-policy/regulations/45-cfr-46/index.html});
the Menlo Report extends an ethical analysis to ICT
research in which data, operators, networks, and bystanders create distinctive
risks \citep{menlo2012}. These sources do not determine this study's regulatory
classification. The responsible authority must issue the determination.

Before any participant activity, the final protocol must record approval,
exemption, or other determination; recruitment and consent materials; any waiver
basis; compensation; inclusion and exclusion criteria; withdrawal; data use and
retention; incident response; and the responsible investigator. Consent must
explain the experimental status, ordinary-review fallback, foreseeable identity
and employment risks, data recipients, automated or agent functions, limits of
withdrawal after aggregation, and whom to contact outside the supervisory chain.

The following anti-coercion constraints are binding design requirements:

\begin{itemize}
  \item recruitment and consent occur outside direct supervisory chains wherever
  possible;
  \item declining, withdrawing, appealing, challenging, or receiving an adverse
  private outcome cannot change employment, funding, authorship, publication,
  ordinary-review access, or standing outside the pilot;
  \item compensation is not contingent on agreement, favorable review, continued
  participation, or forgoing an appeal;
  \item managers cannot request persona linkage, identity opening, assignment
  overrides, or private histories outside the approved process;
  \item a nonmanager channel accepts questions, withdrawal, complaints, and
  adverse-event reports; and
  \item an independent appeal operator can correct process errors and append
  provenance-preserving corrections without erasing the original record.
\end{itemize}

The investigators will publish a positionality and control statement before
recruitment. It will disclose who designed the protocol and tooling; who employs,
funds, supervises, or evaluates participants; who controls the domain, repository,
compute, identity, incident process, and prospective publication venue; and who
benefits from a positive conclusion. Governance meetings and qualitative sessions
will preserve role and status differences. Apparent consensus is not treated as
independent agreement, and a domain's reasoned rejection of \tdrg{} is an
admissible result.

\begin{evidencegate}[Enforceability gate]
Written protections are insufficient if management or operators can bypass them.
Before Stage~1, an independent reviewer must exercise withdrawal, appeal, opening,
rollback, and ordinary-review fallback against the deployed configuration. Failure
or conflicted custody is a no-go condition.
\end{evidencegate}

\section{Staged study design}

\subsection{Stage 0: non-participant readiness}

Stage~0 freezes the intervention, comparator description, data-flow map, operator
roles, opening procedure, threat model, instruments, synthetic cases, adjudication
rubric, and analysis code. Synthetic or appropriately sanitized fixtures will test
malformed events, stale evidence, unauthorized bridge import, cross-topic leakage,
policy divergence, fork and exit, duplicate attempts, appeals, operator loss,
identity-linkage attacks, and disclosure suppression. Tabletop exercises involving
human responses are research activity unless the responsible authority determines
otherwise; ordinary software tests must not quietly become participant studies.

Stage~0 stops if the tested Paper~03 release, incident operator, privacy controls,
synthetic reproduction path, or rollback is unavailable. Passing synthetic tests
establishes only named behavior in named fixtures.

\subsection{Stage 1: consented shadow pilot}

Stage~1 may start only after the full human-subject and preregistration gates. Real
or realistic review cases are used only under approved access and consent or waiver
terms. The intervention runs in parallel with the ordinary process. Experimental
routing recommendations, personas, standing, and explanations are sequestered from
official decision makers until the ordinary disposition is fixed, except where an
approved safety procedure requires narrowly scoped disclosure.

The stage begins with a limited sentinel cohort and proceeds in prespecified
batches. After each batch, an independent safety group reviews only the information
needed for privacy, coercion, burden, exclusion, and protocol-fidelity decisions.
Investigators cannot expand enrollment in response to favorable effectiveness
signals. Any progression pause, amendment, or termination is logged before the
next batch.

\subsection{Stage 2: bounded operational pilot}

Stage~2 is a different intervention and requires a separate ethics confirmation or
amendment, domain consent, threat review, and preregistration. Only after Stage~1
meets all safety and measurement criteria may \tdrg{} receive a narrowly bounded
operational role with rollback and ordinary-review fallback. No reputation output
may be used for public ranking, employment, funding, discipline, or automatic
publication decisions.

\section{Outcomes and measurement plan}

Table~\ref{tab:outcomes} is a candidate measurement catalog, not a commitment to a
composite success score. Co-primary, secondary, exploratory, and safety outcomes
will be frozen before recruitment. Severe events are reported individually under
protected access and cannot be averaged away by review yield or satisfaction.

\begin{longtable}{p{0.17\linewidth}p{0.38\linewidth}p{0.35\linewidth}}
\caption{Candidate outcome families and interpretation constraints.}\label{tab:outcomes}\\
\toprule
Family & Candidate measures & Required interpretation constraint \\
\midrule
\endfirsthead
\toprule
Family & Candidate measures & Required interpretation constraint \\
\midrule
\endhead
Review process & Adjudicated material defects and evidence gaps; false alarms;
unresolved findings; later reversals; explanation completeness; reviewer-hours &
Retain adjudication uncertainty; disaggregate by domain, topic, risk, and
experience; agreement is not truth. \\
Timeliness & Invitation, assignment, review, appeal, and revision latency;
abandonment; backlog & Separate initial learning cost from later batches; report
censoring and withdrawals. \\
Routing and capture & Assignment concentration; root and bridge dependence; policy
overrides; unauthorized-change attempts; staged-fault radius & Condition on the
typed control graph; do not infer independence from graph distance alone. \\
Minority and newcomer & Dissent survival; premature adverse updates; waiting time;
supervision; false exclusion; progression; exit & Protect small specialties and
vulnerable groups from disclosure and stigmatizing labels. \\
Privacy and accountability & Effective eligible pool; linkage precision/recall;
suspected and confirmed reidentification; duplicate attempts; appeal/opening
frequency and correctness & Report by adversary and operator assumptions; suppress
small cells; absence of observed incidents is not safety. \\
Understanding and legitimacy & Comprehension tasks; explanation usability;
perceived procedural fairness; willingness to challenge, appeal, withdraw, or exit
& Self-report is not proof of correctness or protection; preserve role and power
differences. \\
Labor and resources & Human and agent hours by role; compensation; governance,
security, and support time; compute; storage; service dependencies & Report who
bears each cost and who receives benefit, including unpaid and hidden labor. \\
Adverse effects & Retaliation; coercion; chilling; exclusion; harassment;
confidentiality breach; false allegation; standing misuse; automation overreliance
& Do not collapse severity into a mean or trade a severe event against average
benefit. \\
\bottomrule
\end{longtable}

Adjudication will use preconstructed or independently selected cases with known,
sanitized defects where feasible and a panel blinded to experimental standing and
participant identity. The panel will classify each finding as material, minor,
false alarm, unresolved, or outside scope and record disagreement. For live cases,
later author or editor agreement is evidence about process response, not ground
truth. Claims whose correctness remains scientifically disputed stay unresolved.

Qualitative work, if approved, will use purposive sampling across domains, roles,
seniority, employment relationships, experience, dissenting positions, and exit
status. The analysis plan will freeze the interview guide, coding approach, audit
trail, reflexivity statement, negative-case search, disagreement process, and
integration with quantitative evidence. Quotes or paraphrases will undergo
disclosure review for rare expertise, specialist facts, and stylometry.

\section{Anonymity, privacy, and data governance}

The threat model will treat authors, reviewers, domain members, line managers,
investigators, platform operators, credential issuers, publishers, analysts,
trustees, outsiders, and coalitions as distinct adversaries. For each adversary it
will state accessible content and metadata, timing, auxiliary knowledge, attack
objective, authority, mitigations, tests, and residual risk.

Minimum attack families are: timing and assignment correlation; rare expertise or
attribute disclosure; prose stylometry and specialist facts; relationship-graph
inference; cross-case linkage; operator collusion; small-cell inference; credential
or opening abuse; legal or safety demands; and linkage through external services.
Candidate mitigations include batching, delays, access separation, coarse
credentials, minimum effective pools, fresh case keys, disclosure-limited
explanations, separation of duties, threshold opening, immutable access logs,
independent audit, suppression, aggregation, and controlled results. Each mitigation
has operational costs and may fail.

Civil identity and private trust histories will not enter the public research
dataset. The identity map, if one is necessary, is held by the minimum set of
approved operators under a separate retention and destruction schedule. Research
identifiers do not double as review personas. Public artifacts use synthetic
identities and disclosure-reviewed aggregates; controlled artifacts use purpose-
limited access, logging, and expiry. Investigators will document repository and
compute jurisdictions, backups, model providers, telemetry, and deletion limits.

An authorized opening must have a predeclared purpose, threshold, decision maker,
minimum disclosure, logging, appeal and correction path, and participant notice
where legally and ethically permitted. Emergency access does not retroactively
validate an anonymity claim. Suspected duplicate participation remains a process
event until adjudicated and cannot automatically become a public conduct finding.

\section{Analysis plan}

Before recruitment, the statistical analysis plan will freeze co-primary
estimands, smallest meaningful differences or precision targets, sample-size
rationale, allocation, analysis populations, clustering, repeated observations,
period and learning effects, covariates, strata, interactions, multiplicity,
missingness, withdrawal, censoring, late adjudication, and protocol deviations.
Effect estimates will be accompanied by uncertainty intervals; a nonsignificant
test will not be described as equivalence unless an equivalence or noninferiority
design and margin were preregistered.

The default analysis hierarchy is case within domain and period, with reviewer
cross-classification where assignments recur. Domain-specific estimates are
reported before any pooled estimate. A pooled model, if used, must state its
exchangeability assumptions and show sensitivity to excluding each domain and to
reclassifying common control. Outcome-dependent continuation, optional stopping,
and data-driven subgroup construction are prohibited for confirmatory claims.

Missing data are typed by mechanism: participant withdrawal, access restriction,
operator loss, incomplete review, censored appeal, unresolved adjudication,
instrument failure, and disclosure suppression. Withdrawal does not become a
negative performance outcome. The preregistration will state which information may
be retained to honor legal, scientific-integrity, or safety obligations and which
is destroyed; analysis will report the resulting selection limitations.

No algorithmic reviewer-quality label is an outcome. Where a local standing view
changes, analysis will decompose the triggering typed events and policy rules. The
analysis will not use editor agreement, majority vote, author status, citation
count, rank, or institutional prestige as a correctness label. Agent-produced
annotations remain separate from human adjudication, and their model, prompt,
tools, context, and human override are recorded as part of the resource and
provenance package.

\subsection{Progression analysis}

Stage progression is conjunctive rather than compensatory. Measurement fidelity,
privacy, retaliation, exclusion, burden, governance integrity, rollback, and appeal
must each remain within preregistered bounds. A favorable review-yield estimate
cannot offset a failed safety bound. The independent safety group will return
\texttt{continue}, \texttt{pause/remediate}, or \texttt{stop}; investigators will
not receive unblinded effectiveness summaries before a progression decision unless
the preregistration establishes a scientific necessity.

\section{Adverse events, appeals, and stop rules}

Any participant, operator, domain authority, investigator, or independent safety
member may invoke a protected pause. The affected function pauses immediately upon:

\begin{itemize}
  \item unauthorized identity linkage, opening, or private-history access;
  \item credible retaliation, coercion, harassment, employment or funding threat,
  or publication penalty;
  \item material confidentiality or restricted-data breach;
  \item experimental output affecting official disposition outside approved scope;
  \item automatic sanction based on a vote, graph anomaly, disagreement, or
  allegation;
  \item bypass of assignment, appeal, opening, or ordinary-review fallback;
  \item burden, distress, exclusion, or uncompensated labor beyond the registered
  bound;
  \item defeat of the declared privacy or policy boundary by an ordinary operator;
  \item absence or conflict of the incident or appeal operator; or
  \item protocol or tooling drift that invalidates consent, threat assumptions, or
  preregistration.
\end{itemize}

The final safety plan will identify the pause recipient, minimum information needed
for triage, protection against retaliation, investigation process, opening limits,
remediation, participant communication, regulatory reporting, restart authority,
and aggregate reporting rule. The appeal record preserves the original event and
appends a correction, scope change, or reasoned denial. Scientific disagreement is
not misconduct. A safety stop is a scientifically reportable outcome, not a failed
result to be hidden.

\section{Disclosure-safe reproducibility}

Reproducibility does not override consent, confidentiality, data-use agreements,
law, or institutional obligations. The public or broadly shared package will
include the protocol, preregistration, amendments, releasable consent/waiver basis,
instrument specifications, code, schemas, environment, configuration, randomness,
compute description, synthetic fixtures, analysis logs, disclosure rules, and an
immutable manifest. Every table and figure will map to a claim identifier, analysis
run, input class, and disclosure decision.

For restricted participant data, the package will publish variable definitions,
transformations, missingness rules, aggregate consistency checks, and a synthetic
dataset that exercises every analysis branch without representing a participant.
Where allowed, an independent analyst will rerun the locked analysis inside the
controlled environment and sign an attestation naming the code digest, data digest
or custodian-bound identifier, environment, outputs, deviations, and disclosure
class. An attestation establishes execution under stated custody; it does not make
the raw data public or prove scientific validity.

The resource manifest is multiaxis. For every reproducibility or review function it
records human expertise, agent capability and supervision, governance and support
labor, compute, accelerator or scheduler requirements, storage, restricted access,
external services, expected duration, and freshness. A disclosure-safe downselect
may permit digest audit, bounded recomputation, or provenance-only assessment when
full rerun is infeasible, but the resulting evidence level and excluded claims are
explicit. A digest of a prohibited input is not a substitute for independent
scientific access.

Later reproduction failures create append-only events linked to the affected
claims. The record states whether failure arose from unavailable data, expired
authority, environment drift, dependency change, compute loss, undocumented
operator action, or scientific mismatch. Claims may be narrowed or superseded;
the original record is not silently overwritten.

\section{Evidence promotion and planned reporting}

This draft contains design choices, background claims, and open hypotheses. It
contains no implementation or human-subject evidence. A future empirical sentence
may enter the Results section only when its ledger entry names the estimand,
population, comparison, uncertainty, data version, analysis run, deviations,
adversary where relevant, and disclosure class. Qualitative claims must retain
negative cases and positionality while protecting participants. ``No incidents
observed'' cannot be promoted to ``safe.'' A failed linkage attack cannot be
promoted to ``anonymous'' outside its adversary, sample, auxiliary information,
and observation window.

The future Results architecture is locked as follows, but no result prose may be
written before data lock:

\begin{enumerate}
  \item participant and case flow, including refusal and withdrawal;
  \item implementation fidelity and every protocol deviation;
  \item all co-primary outcomes with uncertainty;
  \item domain, role, and preregistered subgroup heterogeneity;
  \item qualitative themes with negative cases and disagreement;
  \item privacy attacks, appeals, openings, and adverse events;
  \item labor and multiaxis resource burden; and
  \item the preregistered go/no-go decision.
\end{enumerate}

Null effects, conflicting domain outcomes, excessive cost, failed measurement, and
a no-go safety conclusion are valid reportable outcomes. No result may be reframed
as success merely because the system executed as designed.

\section{Limitations and transfer conditions}

The pilot would study a small number of self-selected scientific communities using
a specific protocol, release, comparator, and institutional arrangement. It cannot
establish universal reviewer behavior, field-wide fairness, durable anonymity,
scientific validity, or resistance to adversaries outside the threat model. Shadow
use may understate harms or workload that appear when outputs acquire operational
consequences. Conversely, participant attention and investigator support may
overstate routine performance. Small specialties make privacy and statistical
inference especially difficult.

Domain independence may be only partial. Shared employment, funding, hosting,
identity, legal jurisdiction, or personnel can couple nominally separate domains.
The control graph makes these dependencies inspectable but does not remove them.
Plural views can preserve disagreement while also reproducing local exclusion or
allowing strategic forking. Agent assistance may reduce some labor while creating
new dependence on opaque models, service telemetry, prompt context, and human
oversight.

A larger or operational study is warranted only if outcomes are measurable without
equating agreement with truth; participants can exercise withdrawal, appeal, and
exit; independent domains retain credible control and fork paths; privacy,
retaliation, exclusion, burden, and governance harms remain within each registered
bound; and the full resource package is credible. Any unmet condition yields pause,
redesign, or no-go.

\section{Current gate status}

\begin{tabularx}{\linewidth}{p{0.19\linewidth}X p{0.20\linewidth}}
\toprule
Gate & Requirement & Current state \\
\midrule
G0 Argument & Bounded feasibility/safety protocol; no superiority narrative &
Drafted; independent challenge pending \\
G1 Scholarship & Primary peer-review, power, governance, ethics, privacy, and
preregistration sources & Initial verified set; systematic gap audit pending \\
G2 Methods & Statistics, qualitative, domain, privacy/security, and ethics review &
Not started \\
G3 Artifact & Frozen intervention, deployment manifest, synthetic reproduction,
protected-data controls & Not available \\
G4 Independent review & Reviewer outside protocol authors and participating
management challenges harms and interpretation & Not started \\
G5 Human release & Human authors, investigators, ethics/privacy officials, and
domain authorities approve dissemination as applicable & Not authorized \\
\bottomrule
\end{tabularx}

\section{Conclusion}

A machine-verifiable review record can support inspection, but scaling it into a
peer-review corpus introduces independent actors, plural authority, labor,
coercion, privacy, and institutional capture. The proposed pilot therefore begins
in shadow mode and treats governance and safety as measured outcomes rather than
properties inferred from code. Its central commitment is negative as much as
positive: no global reviewer score, no automatic sanction, no unqualified anonymity
claim, no management bypass, no human-subject work before independent determination,
and no field conclusion before preregistered evidence. Whether localized trust
domains improve any part of scientific review remains an open empirical question.

\appendix
\section{Pre-enrollment gate checklist}

The accountable investigator must attach immutable identifiers for every item:

\begin{enumerate}
  \item written ethics/IRB determination and approved protocol;
  \item recruitment, consent or waiver basis, compensation, and exit procedure;
  \item privacy/security assessment, adversary matrix, and breach response;
  \item domain agreements, typed control graphs, operator conflicts, and fork test;
  \item data-management, retention, jurisdiction, and disclosure plan;
  \item preregistration with outcomes, estimands, precision, allocation, analysis,
  missingness, stop rules, and deviations;
  \item tested protocol/tool release and deployment manifest;
  \item independent adjudication, appeal, incident, and safety operators;
  \item synthetic end-to-end reproduction and rollback evidence; and
  \item human authorization for authorship, AI disclosure, venue, licensing, and
  any release.
\end{enumerate}

If one item is absent or stale, the protocol remains at Stage~0.

\section{Adversary-record template}

For each adversary or coalition, record: identifier; role; authority; accessible
content and metadata; timing; auxiliary data; attack objective; attack procedure;
assumed operator honesty; mitigation; residual risk; quantitative threshold;
observation window; result; incident linkage; disclosure class; and the exact
language permitted in any privacy claim. Coalition tests must include at least line
management plus platform operator, investigator plus analyst, and credential issuer
plus identity/opening operator where those roles are separable.

\section{Evidence-promotion template}

Each proposed empirical statement must include: claim ID; exact text; label
(implementation observation, human-subject evidence, inference, or implication);
estimand; unit; population; comparison; data/custodian identifier; code and
environment digests; run ID; uncertainty; deviations; adjudication status;
adversary and observation window when applicable; subgroup and disclosure class;
support locator; falsifier; supersession relation; and required human sign-off.

\bibliographystyle{plainnat}
\bibliography{references}

\end{document}
